\documentclass[oneside,12pt,a4paper]{article}

\usepackage{amsmath}

\title{SME0340 -- Primeira Prova}
\date{}
\author{}

\usepackage[T1]{fontenc}
\usepackage[utf8]{inputenc}
\usepackage{mathpazo}
\usepackage[brazil]{babel}
\usepackage{microtype}


\begin{document}

    \maketitle

    \section*{Instruções}

    \vspace*{0.5em}
    \textbf{\textsc{Leia atentamente as instruções abaixo antes de começar.}}

    \subsection*{Entrega}

    A prova deverá ser entregue através do sistema eDisciplinas até às 23:59 do dia 02/06/2026. A entrega deve ser através de um único arquivo no formato PDF. O arquivo PDF pode ser obtido através de digitalização de manuscritos ou geração direta (através de MarkDown, LaTeX, softwares de notas manuscritas, et.) desde que seja legível e que seja compatível com todos visualizadores.

    \subsection*{Conteúdo}

    A resposta a cada questão deverá conter:
    \begin{itemize}
        \item Um cabeçalho indicando claramente qual questão está sendo resolvida, incluindo o enunciado da questão (o LaTeX da prova está disponibilizado para ajudar com isso);

        \item O desenvolvimento completo da questão, usando suas palavras para explicar cada um dos passos. Caso seja utilizado algum software para a resolução, explique para que parte e como foi utilizado;

        \item A resposta final, claramente indicada.
    \end{itemize}

    \subsection*{Resolução}

    Você \textbf{poderá} realizar consultas bibliográficas, utilizar software de computação simbólica ou consultas a sistemas de inteligência artificial para resolver as questões das provas. Você \textbf{não poderá} consultar as respostas de colegas. Tire dúvidas com o professor em vez de perguntar para colegas.

    \pagebreak

    \section*{Questão 1}

    Considere a equação diferencial ordinária dada por
    $$
        \frac{\mathrm dy}{\mathrm dx} = e^{-x^2}.
    $$
    Sem resolver a EDO, responda os itens abaixo:
    \begin{enumerate}
        \item Explique por que a solução da ED necessariamente é uma função crescente em qualquer intervalo do eixo $x$.

        \item Quais são os limites $\displaystyle\lim_{x\to-\infty}\mathrm dy/\mathrm dx$ e $\displaystyle\lim_{x\to\infty}\mathrm dy/\mathrm dx$? O que isto implica a respeito da curva da solução quando $x\to\pm\infty$?

        \item Determine um intervalo no qual a solução $y=\phi(x)$ tem concavidade para baixo e um intervalo no qual a curva tem concavidade para cima.
    \end{enumerate}

    \section*{Questão 2}

    \begin{enumerate}
        \item Verifique que $y=-1/(x+c)$ é uma família a um parâmetro de soluções da equação diferencial $y'=y^2$.

        \item Ache uma solução da família no item 1 que satisfaça $y(0)=1$. Então, ache uma solução da família no item 1 que satisfaça $y(0)=-1$. Determine o maior intervalo $I$ de definição da solução de cada problema de valor inicial.

        \item Determine o maior intervalo $I$ de definição para o problema de valor inicial $y'=y^2$, $y(0)=0$.
    \end{enumerate}

    \section*{Questão 3}

    Em cada caso abaixo, ache a solução geral da equação diferencial dada.
    \begin{enumerate}
        \item
        $$
            y\ln(x)\frac{\mathrm dy}{\mathrm dx}=\left(\frac{y+1}{\sqrt x}\right)^2.
        $$
        \item
        $$
            \frac{\mathrm dy}{\mathrm dx}=\frac{xy+3x-y-3}{xy-2x+4y-8}.
        $$
    \end{enumerate}

    \section*{Questão 4}

    \begin{enumerate}
        \item Encontre a solução geral da seguinte EDO:
        $$
            (4xy+3x^2)+(2y+2x^2)\frac{\mathrm dy}{\mathrm dx}=0.
        $$

        \item Mostre que as condições iniciais $y(0)=-2$ e $y(1)=1$ determinam a mesma solução implícita.

        \item Ache soluções explícitas $y_1(x)$ e $y_2(x)$ da equação diferencial do item 1 tal que $y_1(0)=-2$ e $y_2(1)=1$. 
    \end{enumerate}

    \section*{Questão 5}
    Um marca-passo cardíaco consiste em uma chave, uma bateria de tensão constante $E_0$, um capacitor com capacitância constante $C$, e o coração, agindo como um resistor com resistência constante $R$. Quando a chave é fechada, o capacitor se carrega; quando a chave é aberta, o capacitor se descarrega, enviando um estímulo elétrico para o coração. Durante o tempo em que  o coração é estimulado, a tensão $E$ em todo o coração satisfaz a equação diferencial linear
    $$
        \frac{dE}{dt}=-\frac{1}{RC}E
    $$
    Resolva a EDO acima e o PVI associado com condição inicial $E(4)=E_0$.

    \section*{Questão 6}

    O peso $p(t)$ para uma espécie de peixe é dado pela seguinte equação obtida experimentalmente:
    $$
        \frac{dp}{dt}=\alpha p^\frac{2}{3}-\beta p
    $$
    onde $\alpha$ é a constante de anabolismo (taxa de síntese de massa por unidade de superfície do animal) e $\beta$ é a constante de catabolismo (taxa de diminuição de massa por unidade de superfície). Resolva a ED levando em conta que ela é uma equação de Bernoulli. Assumindo que $p(0)=0$, determine em qual instante de tempo o peixe atinge o peso de $\frac{1}{16}\left(\frac{\alpha}{\beta}\right)^3$.

\end{document}
